\section{Reproducibility of the Experiments}
\label{env:reproducibility}

Reproducibility is a central methodological requirement for a simulation-based experimental thesis, especially when the contribution consists of updating a complex toolchain and re-implementing logic that has already been published in the literature. In this context, reproducibility is not limited to the possibility of obtaining a qualitatively similar result, but requires configurations, parameters, software versions, and execution pipelines to be tracked well enough to enable a reliable comparison among experiments.

For OMNeT++/SUMO co-simulation, the main reproducibility issues are:
\begin{itemize}
\item version differences among frameworks and dependencies;
\item undocumented changes in configuration files, such as \texttt{omnetpp.ini}, NED, and SUMO scenarios;
\item inconsistent management of random seeds and the number of replicas;
\item code modifications that are not isolated or not tracked;
\item manual execution and post-processing pipelines that are difficult to replicate.
\end{itemize}
For this reason, the present work adopts an explicit distinction between:
\begin{itemize}
\item the development environment, described in this chapter, namely the tools and roles of the components;
\item the concrete experimental environment, described in the next chapter, namely project organization, build process, configurations, and operational validation.
\end{itemize}

From a practical point of view, an effective reproducibility strategy for this type of thesis includes at least the following elements:
\begin{itemize}
\item precise identification of the versions of the frameworks used;
\item separation between modified source code and experiment configuration files;
\item definition of scenarios and parameters in documented parametric form;
\item automation of repeated executions and output collection;
\item storage of raw metrics together with run metadata, such as seed, scenario, timestamp, and configuration.
\end{itemize}
These aspects are not mere implementation details: they directly affect the credibility of the results and the possibility of correctly interpreting any differences with respect to the reference paper.

Finally, reproducibility also has a specific scientific meaning in this thesis. Since the work consists of a re-implementation and update of a multi-technology system with newer framework versions and with 5G instead of LTE, accurate documentation of the environment and the workflow is an integral part of the contribution: it makes it possible to distinguish what depends on the logic under study, for example \emph{SafeSwitch}, from what depends on the simulation infrastructure and its configurations.
